% GENERATED FILE. Edit canonical lessons, then run npm run generate:books.
% canonical-source-hash: fnv1a64:37172380facbe3fe
% canonical-lessons: ZH-C104-benzi, ZH-W104-ben, ZH-C104-bi, ZH-C104-qianbi, ZH-C104-zhi, ZH-C104-cidian, ZH-R104-review-99-104

\chapter{On the Desk}
\label{ch:zh-on-the-desk}
\hlchaptermodality{104}

\begin{chapteropening}
\textbf{By the end of this chapter:} \emph{I can name a notebook, a pen, a pencil, paper and a dictionary.}

Complete the last lesson of chapter 104: I can name a notebook, a pen, a pencil, paper and a dictionary.
\end{chapteropening}

\section[\texorpdfstring{\zh{本子} (běnzi)}{本子 (běnzi)}]{\zh{本子} (běnzi) — notebook}

\label{lesson:ZH-C104-benzi}



\begin{quote}
Before the new one: say the Mandarin for to repair, then the Mandarin for broken.
\end{quote}

\subsection*{You'll want to know: \zh{本子}}
\textbf{\zh{本子}} — \emph{běnzi} — \textquotedblleft{}notebook\textquotedblright{}.

\begin{quote}
\textbf{\zh{我能用你的本子吗}} \emph{wǒ néng yòng nǐ de běnzi ma} — Could I use your notebook?
\end{quote}

The new character, in pieces: \zh{本} — the tree piece of \zh{机}, standing alone, with a short bar across its foot. The next lesson puts it in your hand.

\begin{cousinweb}
\textbf{\zh{本子}} (\emph{běnzi}) is built from \zh{本} and the noun ending \zh{子}, both inherited from Middle Chinese, \emph{pwonX} and \emph{tsiX} (Guangyun, in Baxter's transcription). The dictionary glosses the first 'root' and explains its shape as a tree with a mark at its foot; how it came to mean a book is not settled.
\end{cousinweb}

\begin{grammarlens}[title={What you've built}]
One more word for everyday things.
\end{grammarlens}

\subsection*{Your turn}
\begin{itemize}
\raggedright
  \item \emph{Say it:} \emph{běnzi}
  \item \emph{Say it:} \emph{běnzi}, once more
  \item \emph{Say it:} \emph{wǒ néng yòng nǐ de běnzi ma} — Could I use your notebook?
  \item \emph{Recall:} say the Mandarin for to repair, then the Mandarin for broken, then say \emph{běnzi} again
\end{itemize}

\begin{tcolorbox}[breakable,colback=teal!4,colframe=teal!35!black,title={Before you move on}]
What does \zh{本子} mean? (\textquotedblleft{}Notebook\textquotedblright{}.) Say it once more.
\end{tcolorbox}
\section[\texorpdfstring{\zh{本} (běn)}{本 (běn)}]{\zh{本} — the tree piece, and a bar at its foot}

\label{lesson:ZH-W104-ben}



\begin{quote}
The word first, then the character inside it.

\begin{itemize}
\raggedright
  \item \emph{Recall:} say \emph{běnzi} — \textbf{R1}, one lesson back, before it decays
\end{itemize}

\begin{quote}
\zh{机} \zh{子}
\end{quote}

\emph{Read it:} each one aloud, with nothing said first — \emph{jī}, \emph{zǐ}
\end{quote}

\subsection*{Script — the tree piece, and a bar at its foot}
\hlblockfigure{\detokenize{figures/ZH-W104-ben-filmstrip.pdf}}{How \zh{本} is written, stroke by stroke}

\begin{quote}
\zh{本}
\end{quote}

\textbf{The tree piece, four strokes}, as on the left of \zh{机} but full width: a horizontal, a vertical down through it, a left-falling stroke and a long right-falling one.

\textbf{A short horizontal} across the vertical, low down: the mark at the root.

4 + 1 = five.

\subsection*{Your turn}
\begin{itemize}
\raggedright
  \item \emph{Write it:} \zh{本}, all five, with the model in front of you
  \item \emph{Cover:} write it again from memory, and count the pieces, not the strokes
\end{itemize}

\begin{tcolorbox}[breakable,colback=teal!4,colframe=teal!35!black,title={Before you move on}]
How many strokes in \zh{本}? (\textbf{Five}.) Where is the short bar? (\textbf{Low down, across the vertical}.) Which character you know has the same tree piece? (\textbf{\zh{机}}.)

Source: \href{https://raw.githubusercontent.com/chanind/hanzi-writer-data/68d10a4b21150cae5e1ebbd223eed289cf32d90c/data/\%E6\%9C\%AC.json}{Hanzi Writer Data \zh{本}, pinned stroke paths and medians}.
\end{tcolorbox}
\section[bǐ]{bǐ — pen}

\label{lesson:ZH-C104-bi}



\begin{quote}
Before the new one: say the Mandarin for broken, then the Mandarin for notebook.
\end{quote}

\subsection*{You'll want to know: bǐ}
\emph{bǐ} — \textquotedblleft{}pen\textquotedblright{}.

\begin{quote}
\emph{nǐ yǒu bǐ ma} — Have you got a pen?
\end{quote}

Its characters come later; for now it is a word for the ear.

\begin{cousinweb}
\emph{bǐ} is inherited from Middle Chinese \emph{pit} (Guangyun, in Baxter's transcription), the writing brush; the dictionary repeats a story that a general of the Qin dynasty made the first one.
\end{cousinweb}

\begin{grammarlens}[title={What you've built}]
One more word for everyday things.
\end{grammarlens}

\subsection*{Your turn}
\begin{itemize}
\raggedright
  \item \emph{Say it:} \emph{bǐ}
  \item \emph{Say it:} \emph{bǐ}, once more
  \item \emph{Say it:} \emph{nǐ yǒu bǐ ma} — Have you got a pen?
  \item \emph{Recall:} say the Mandarin for broken, then the Mandarin for notebook, then say \emph{bǐ} again
\end{itemize}

\begin{tcolorbox}[breakable,colback=teal!4,colframe=teal!35!black,title={Before you move on}]
What does bǐ mean? (\textquotedblleft{}Pen\textquotedblright{}.) Say it once more.
\end{tcolorbox}
\section[qiānbǐ]{qiānbǐ — pencil}

\label{lesson:ZH-C104-qianbi}



\begin{quote}
Before the new one: say the Mandarin for notebook, then the Mandarin for pen.
\end{quote}

\subsection*{You'll want to know: qiānbǐ}
\emph{qiānbǐ} — \textquotedblleft{}pencil\textquotedblright{}.

\begin{quote}
\emph{wǒ yòng qiānbǐ xiě zì} — I write with a pencil.
\end{quote}

Its characters come later; for now it is a word for the ear.

\begin{cousinweb}
\emph{qiānbǐ} is built from qiān 'lead' and bǐ 'pen', both inherited from Middle Chinese, \emph{ywen} and \emph{pit} (Guangyun, in Baxter's transcription); the dictionary glosses the first as a grey-blue metal. The regular change from \emph{ywen} gives yuán, not qiān; how the word came to its sound is not settled.
\end{cousinweb}

\begin{grammarlens}[title={What you've built}]
One more word for everyday things.
\end{grammarlens}

\subsection*{Your turn}
\begin{itemize}
\raggedright
  \item \emph{Say it:} \emph{qiānbǐ}
  \item \emph{Say it:} \emph{qiānbǐ}, once more
  \item \emph{Say it:} \emph{wǒ yòng qiānbǐ xiě zì} — I write with a pencil.
  \item \emph{Recall:} say the Mandarin for notebook, then the Mandarin for pen, then say \emph{qiānbǐ} again
\end{itemize}

\begin{tcolorbox}[breakable,colback=teal!4,colframe=teal!35!black,title={Before you move on}]
What does qiānbǐ mean? (\textquotedblleft{}Pencil\textquotedblright{}.) Say it once more.
\end{tcolorbox}
\section[zhǐ]{zhǐ — paper}

\label{lesson:ZH-C104-zhi}



\begin{quote}
Before the new one: say the Mandarin for pen, then the Mandarin for pencil.
\end{quote}

\subsection*{You'll want to know: zhǐ}
\emph{zhǐ} — \textquotedblleft{}paper\textquotedblright{}.

\begin{quote}
\emph{qǐng gěi wǒ yì zhāng zhǐ} — Could I have a sheet of paper, please?
\end{quote}

Its characters come later; for now it is a word for the ear.

\begin{cousinweb}
\emph{zhǐ} is inherited from Middle Chinese \emph{tsyeX} (Guangyun, in Baxter's transcription), where the dictionary says paper is as smooth as a whetstone, and that in the Later Han it was made from old fishing nets and tree bark.
\end{cousinweb}

\begin{grammarlens}[title={What you've built}]
One more word for everyday things.
\end{grammarlens}

\subsection*{Your turn}
\begin{itemize}
\raggedright
  \item \emph{Say it:} \emph{zhǐ}
  \item \emph{Say it:} \emph{zhǐ}, once more
  \item \emph{Say it:} \emph{qǐng gěi wǒ yì zhāng zhǐ} — Could I have a sheet of paper, please?
  \item \emph{Recall:} say the Mandarin for pen, then the Mandarin for pencil, then say \emph{zhǐ} again
\end{itemize}

\begin{tcolorbox}[breakable,colback=teal!4,colframe=teal!35!black,title={Before you move on}]
What does zhǐ mean? (\textquotedblleft{}Paper\textquotedblright{}.) Say it once more.
\end{tcolorbox}
\section[cídiǎn]{cídiǎn — dictionary}

\label{lesson:ZH-C104-cidian}



\begin{quote}
Before the new one: say the Mandarin for pencil, then the Mandarin for paper.
\end{quote}

\subsection*{You'll want to know: cídiǎn}
\emph{cídiǎn} — \textquotedblleft{}dictionary\textquotedblright{}.

\begin{quote}
\emph{wǒ yǒu yì běn hànyǔ cídiǎn} — I've got a Chinese dictionary.
\end{quote}

Its characters come later; for now it is a word for the ear.

\begin{cousinweb}
\emph{cídiǎn} is built from cí 'words' and diǎn 'a standard work', both inherited from Middle Chinese, \emph{zi} and \emph{tenX} (Guangyun, in Baxter's transcription); the dictionary glosses the second as a law and a classic.
\end{cousinweb}

\begin{grammarlens}[title={What you've built}]
One more word for everyday things.
\end{grammarlens}

\subsection*{Your turn}
\begin{itemize}
\raggedright
  \item \emph{Say it:} \emph{cídiǎn}
  \item \emph{Say it:} \emph{cídiǎn}, once more
  \item \emph{Say it:} \emph{wǒ yǒu yì běn hànyǔ cídiǎn} — I've got a Chinese dictionary.
  \item \emph{Recall:} say the Mandarin for pencil, then the Mandarin for paper, then say \emph{cídiǎn} again
\end{itemize}

\begin{tcolorbox}[breakable,colback=teal!4,colframe=teal!35!black,title={Before you move on}]
What does cídiǎn mean? (\textquotedblleft{}Dictionary\textquotedblright{}.) Say it once more.
\end{tcolorbox}
\section[(dialogue)]{Looking back: from Getting Washed to On the Desk}

\label{lesson:ZH-R104-review-99-104}



\begin{quote}
Say the Mandarin for to wash and to brush one's teeth.
\end{quote}

\subsection*{Your turn}
\begin{itemize}
\raggedright
  \item \emph{Say it:} to wash, to brush one's teeth, to wash one's face, to have a shower, to comb one's hair
  \item \emph{Say it:} bank, supermarket, park, library, post office
  \item \emph{Say it:} kitchen, cabbage, chopsticks, plate, knife
  \item \emph{Say it:} soft drink, juice, cola, beer, biscuit
  \item \emph{Say it:} to use, to take hold of, to throw away, to repair, broken
  \item \emph{Say it:} notebook, pen, pencil, paper, dictionary
  \item \emph{Read it:} \zh{洗} \zh{菜} \zh{用} \zh{本} — each one aloud, then a word that holds it
\end{itemize}

\begin{tcolorbox}[breakable,colback=teal!4,colframe=teal!35!black,title={Before you move on}]
To wash? (\textbf{\zh{洗}}, \emph{xǐ}.) To brush one's teeth? (\emph{shuāyá}.) To wash one's face? (\emph{xǐliǎn}.) To have a shower? (\emph{xǐzǎo}.) To comb one's hair? (\emph{shūtóu}.) Bank? (\emph{yínháng}.) Supermarket? (\emph{chāoshì}.) Park? (\emph{gōngyuán}.) Library? (\emph{túshūguǎn}.) Post office? (\emph{yóujú}.) Kitchen? (\emph{chúfáng}.) Cabbage? (\textbf{\zh{白菜}}, \emph{báicài}.) Chopsticks? (\emph{kuàizi}.) Plate? (\emph{pánzi}.) Knife? (\emph{dāo}.) Soft drink? (\emph{yǐnliào}.) Juice? (\emph{guǒzhī}.) Cola? (\emph{kělè}.) Beer? (\emph{píjiǔ}.) Biscuit? (\emph{bǐnggān}.) To use? (\textbf{\zh{用}}, \emph{yòng}.) To take hold of? (\emph{ná}.) To throw away? (\emph{rēng}.) To repair? (\emph{xiū}.) Broken? (\emph{huài}.) Notebook? (\textbf{\zh{本子}}, \emph{běnzi}.) Pen? (\emph{bǐ}.) Pencil? (\emph{qiānbǐ}.) Paper? (\emph{zhǐ}.) Dictionary? (\emph{cídiǎn}.)
\end{tcolorbox}
